\documentclass[11pt]{article}
\pagestyle{empty}

\setlength{\topmargin}{-.75in}
\setlength{\textheight}{9.2in}
\setlength{\oddsidemargin}{-.25in}
\setlength{\evensidemargin}{-.25in}
\setlength{\textwidth}{6.75in}
\setlength{\parskip}{4pt}

\usepackage{amsmath}
\usepackage{amsfonts}
\usepackage{amsthm}
\usepackage{amssymb}
\usepackage{latexsym}
\usepackage{float}
\usepackage{graphicx}
\usepackage{mathtools}
\usepackage{booktabs}
\DeclarePairedDelimiter\ceil{\lceil}{\rceil}
\DeclarePairedDelimiter\floor{\lfloor}{\rfloor}

% Figure paths that work whether the document is compiled from the repo root
% or from within PSET_2.
\newcommand{\SlideSeventeenFigurePath}{slide17_sampling_density_rho1.png}
\newcommand{\SlideEighteenFigurePath}{slide18_sampling_distribution_phat.png}
\newcommand{\QFiveARFigurePath}{question5_output/figures/q5_ar_log_spectrum.pdf}
\newcommand{\QFiveKernelFigurePath}{question5_output/figures/q5_kernel_log_spectrum.pdf}
\IfFileExists{\SlideSeventeenFigurePath}{}{
  \renewcommand{\SlideSeventeenFigurePath}{PSET_2/slide17_sampling_density_rho1.png}
  \renewcommand{\SlideEighteenFigurePath}{PSET_2/slide18_sampling_distribution_phat.png}
  \renewcommand{\QFiveARFigurePath}{PSET_2/question5_output/figures/q5_ar_log_spectrum.pdf}
  \renewcommand{\QFiveKernelFigurePath}{PSET_2/question5_output/figures/q5_kernel_log_spectrum.pdf}
}

% Question 5 generated values
\newcommand{\QFiveObs}{864}
\newcommand{\QFiveLagNSA}{25}
\newcommand{\QFiveLagSA}{5}
\newcommand{\QFiveARAnnualNSA}{-0.18}
\newcommand{\QFiveARAnnualSA}{-3.97}
\newcommand{\QFiveKernelAnnualNSA}{-1.03}
\newcommand{\QFiveKernelAnnualSA}{-5.54}
\newcommand{\QFiveARAnnualGap}{3.79}
\newcommand{\QFiveKernelAnnualGap}{4.51}
\newcommand{\QFiveARContrastNSA}{0.02}
\newcommand{\QFiveARContrastSA}{-0.00}
\newcommand{\QFiveKernelContrastNSA}{0.01}
\newcommand{\QFiveKernelContrastSA}{0.00}
\newcommand{\QFiveKernelBandwidth}{10}

% Question 6 generated values
\newcommand{\QSixRawStart}{1948Q1}
\newcommand{\QSixRawEnd}{2019Q4}
\newcommand{\QSixRawN}{288}
\newcommand{\QSixEstStart}{1948Q4}
\newcommand{\QSixEstEnd}{2019Q3}
\newcommand{\QSixEstN}{284}
\newcommand{\QSixBandwidth}{5}
\newcommand{\QSixCOne}{-0.000159}
\newcommand{\QSixGFOne}{0.8225}
\newcommand{\QSixGBOne}{0.2069}
\newcommand{\QSixLamOne}{-0.000013}
\newcommand{\QSixCOneSE}{0.000378}
\newcommand{\QSixGFOneSE}{0.1657}
\newcommand{\QSixGBOneSE}{0.1257}
\newcommand{\QSixLamOneSE}{0.000063}
\newcommand{\QSixCTwo}{-0.000255}
\newcommand{\QSixGFTwo}{0.8363}
\newcommand{\QSixGBTwo}{0.1937}
\newcommand{\QSixLamTwo}{0.000001}
\newcommand{\QSixCTwoSE}{0.000350}
\newcommand{\QSixGFTwoSE}{0.1308}
\newcommand{\QSixGBTwoSE}{0.0994}
\newcommand{\QSixLamTwoSE}{0.000059}
\newcommand{\QSixJStat}{1.6043}
\newcommand{\QSixJPValue}{0.6584}

\def\im{\mathrm{im}}
\def\H{\mathbf{H}}
\def\qbar{\overline{q}}
\def\SO{\mathrm{SO}}
\def\OO{\mathrm{O}}
\def\ZZ{\mathrm{Z}}
\def\C{\mathbf{C}}
\def\R{\mathbf{R}}
\def\Q{\mathbf{Q}}
\def\Z{\mathbf{Z}}
\def\F{\mathbf{F}}
\def\GL{\mathrm{GL}}
\def\SL{\mathrm{SL}}
\def\PGL{\mathrm{PGL}}
\def\PSL{\mathrm{PSL}}
\def\Aut{\mathrm{Aut}}
\def\Inn{\mathrm{Inn}}
\def\Out{\mathrm{Out}}
\def\xbar{\overline{x}}
\def\ybar{\overline{y}}
\def\RO{\mathrm{RO}}

\title{Advanced Time Series PSET 2}
\author{Dhruv Kohli}
\date{April 2026}

\begin{document}

\maketitle

\section{Problem 1}

Let
\[
y_t=\rho y_{t-1}+u_t,\qquad
u_t=\sigma_t\varepsilon_t,\qquad
\sigma_t^2=\nu+\alpha u_{t-1}^2+\beta \sigma_{t-1}^2,
\qquad
\varepsilon_t \overset{i.i.d.}{\sim} N(0,1).
\]
Assume $\{y_t,\sigma_t\}$ is strictly stationary and $\sigma_t>0$. Define the past information set
\[
\mathcal F_{t-1}=\sigma(y_{t-1},y_{t-2},\dots).
\]
Since
\[
u_{t-j}=y_{t-j}-\rho y_{t-j-1}\qquad (j\ge 1),
\]
the past values $\{u_{t-1},u_{t-2},\dots\}$ are measurable with respect to $\mathcal F_{t-1}$. Iterating the GARCH recursion backward therefore shows that $\sigma_t^2$ is also $\mathcal F_{t-1}$-measurable.

\paragraph{(i)}

We first have,
\[
E(u_t^2)=E(\sigma_t^2\varepsilon_t^2)=E(\sigma_t^2)E(\varepsilon_t^2)=E(\sigma_t^2),
\]
since $\varepsilon_t$ is independent of $\sigma_t$ and $E(\varepsilon_t^2)=1$.

Taking expectations in the GARCH recursion,
\[
E(\sigma_t^2)=\nu+\alpha E(u_{t-1}^2)+\beta E(\sigma_{t-1}^2).
\]
By stationarity,
\[
E(\sigma_t^2)=E(\sigma_{t-1}^2)\equiv m,
\qquad
E(u_{t-1}^2)=E(\sigma_{t-1}^2)=m.
\]
Thus,
\[
m=\nu+\alpha m+\beta m,
\]
so
\[
{E(\sigma_t^2)=\frac{\nu}{1-\alpha-\beta}},
\]
provided $\alpha+\beta<1$.

Next, compute the unconditional variance of $y_t$. Since
\[
E(u_t\mid \mathcal F_{t-1})=E(\sigma_t\varepsilon_t\mid \mathcal F_{t-1})=\sigma_t E(\varepsilon_t)=0,
\]
we have
\[
E(y_t)=\rho E(y_{t-1})+E(u_t).
\]
By stationarity, $E(y_t)=E(y_{t-1})$, so
\[
E(y_t)=0.
\]
Therefore,
\begin{align*}
Var(y_t)
&=E(y_t^2)
=E\bigl((\rho y_{t-1}+u_t)^2\bigr) \\
&=\rho^2 E(y_{t-1}^2)+2\rho E(y_{t-1}u_t)+E(u_t^2).
\end{align*}
Now
\[
E(y_{t-1}u_t)=E\!\left[y_{t-1}E(u_t\mid \mathcal F_{t-1})\right]=0,
\]
so
\[
Var(y_t)=\rho^2Var(y_{t-1})+E(u_t^2).
\]
Using stationarity again,
\[
Var(y_t)=Var(y_{t-1})\equiv \gamma_0,
\]
thus,
\[
\gamma_0=\rho^2\gamma_0+\frac{\nu}{1-\alpha-\beta}.
\]
Therefore,
\[
{Var(y_t)=\frac{\nu}{(1-\rho^2)(1-\alpha-\beta)}},
\]
provided $|\rho|<1$ and $\alpha+\beta<1$.

\paragraph{(ii)}

Since
\[
y_t=\rho y_{t-1}+u_t,
\]
we have
\[
E(y_t\mid \mathcal F_{t-1})
=\rho y_{t-1}+E(u_t\mid \mathcal F_{t-1})
=\rho y_{t-1}.
\]
Thus,
\[
{E(y_t\mid y_{t-1},y_{t-2},\dots)=\rho y_{t-1}.}
\]

Also,
\[
Var(y_t\mid \mathcal F_{t-1})
=Var(u_t\mid \mathcal F_{t-1}).
\]
Because $\sigma_t$ is $\mathcal F_{t-1}$-measurable and $\varepsilon_t\sim N(0,1)$,
\[
Var(u_t\mid \mathcal F_{t-1})
=Var(\sigma_t\varepsilon_t\mid \mathcal F_{t-1})
=\sigma_t^2 Var(\varepsilon_t)
=\sigma_t^2.
\]
We then get
\[
{Var(y_t\mid y_{t-1},y_{t-2},\dots)=\sigma_t^2.}
\]

Since $u_{t-1}=y_{t-1}-\rho y_{t-2}$, we may also write
\[
{\sigma_t^2=\nu+\alpha (y_{t-1}-\rho y_{t-2})^2+\beta \sigma_{t-1}^2.}
\]

\paragraph{(iii)}

Conditional on $\mathcal F_{t-1}$,
\[
u_t=\sigma_t\varepsilon_t \sim N(0,\sigma_t^2),
\]
so
\[
y_t\mid \mathcal F_{t-1}\sim N(\rho y_{t-1},\sigma_t^2).
\]
Therefore the conditional density of $y_t$ is
\[
f(y_t\mid \mathcal F_{t-1})
=\frac{1}{\sqrt{2\pi \sigma_t^2}}
\exp\!\left(
-\frac{(y_t-\rho y_{t-1})^2}{2\sigma_t^2}
\right).
\]
Thus, the conditional likelihood is
\[
L(\rho,\nu,\alpha,\beta\mid y_0,y_{-1},\sigma_0^2)
=\prod_{t=1}^T
\frac{1}{\sqrt{2\pi \sigma_t^2}}
\exp\!\left(
-\frac{(y_t-\rho y_{t-1})^2}{2\sigma_t^2}
\right),
\]
where the conditional variances satisfy the recursion
\[
\sigma_t^2=\nu+\alpha (y_{t-1}-\rho y_{t-2})^2+\beta \sigma_{t-1}^2,
\qquad t=1,\dots,T,
\]
with initial condition
\[
{\sigma_0^2=E(\sigma_t^2)=\frac{\nu}{1-\alpha-\beta}.}
\]
In particular,
\[
\sigma_1^2=\nu+\alpha (y_0-\rho y_{-1})^2+\beta \sigma_0^2,
\]
so the likelihood is indeed conditional only on $(y_0,y_{-1})$ and $\sigma_0^2$, exactly as requested in the problem.

Equivalently, the log-likelihood is
\[
\ell(\rho,\nu,\alpha,\beta)
=
-\frac{T}{2}\log(2\pi)
-\frac12\sum_{t=1}^T
\left[
\log \sigma_t^2+\frac{(y_t-\rho y_{t-1})^2}{\sigma_t^2}
\right].
\]

We thus conclude taht,
\[
{E(\sigma_t^2)=\frac{\nu}{1-\alpha-\beta}},\qquad
{Var(y_t)=\frac{\nu}{(1-\rho^2)(1-\alpha-\beta)}},
\]
\[
{E(y_t\mid y_{t-1},y_{t-2},\dots)=\rho y_{t-1}},\qquad
{Var(y_t\mid y_{t-1},y_{t-2},\dots)=\sigma_t^2},
\]
and
\[
{
L(\rho,\nu,\alpha,\beta\mid y_0,y_{-1},\sigma_0^2)
=\prod_{t=1}^T
(2\pi \sigma_t^2)^{-1/2}
\exp\!\left(
-\frac{(y_t-\rho y_{t-1})^2}{2\sigma_t^2}
\right)
}
\]
with
\[
{
\sigma_t^2=\nu+\alpha (y_{t-1}-\rho y_{t-2})^2+\beta \sigma_{t-1}^2,
\qquad
\sigma_0^2=\frac{\nu}{1-\alpha-\beta}.
}
\]

\section{Problem 2}

We use the following design for our simulation:
\[
y_t = 0.6\,y_{t-1} + 0.2\,y_{t-2} + u_t,
\qquad
u_t \overset{i.i.d.}{\sim} N(0,1),
\qquad
T=200,
\qquad
p\in\{1,2,3\},
\]
with \(10{,}000\) Monte Carlo repetitions and lag order chosen by BIC or AIC.

\begin{figure}[H]
\centering
\includegraphics[width=0.72\textwidth]{\SlideSeventeenFigurePath}
\caption{Replication of Slide 17: sampling density of the post-selection estimator \(\hat\rho_1\).}
\end{figure}

\begin{figure}[H]
\centering
\includegraphics[width=0.72\textwidth]{\SlideEighteenFigurePath}
\caption{Replication of Slide 18: sampling distribution of the selected lag order \(\hat p\).}
\end{figure}

The selection frequencies are
\[
P(\hat p_{\mathrm{BIC}}=1,2,3)=(0.3562,\ 0.6291,\ 0.0147),
\]
\[
P(\hat p_{\mathrm{AIC}}=1,2,3)=(0.1046,\ 0.7411,\ 0.1543).
\]
For the post-selection estimator of the first autoregressive coefficient,
\[
E(\hat\rho_{1,\mathrm{BIC}})\approx 0.6195,
\qquad
\mathrm{sd}(\hat\rho_{1,\mathrm{BIC}})\approx 0.0986,
\]
\[
E(\hat\rho_{1,\mathrm{AIC}})\approx 0.5960,
\qquad
\mathrm{sd}(\hat\rho_{1,\mathrm{AIC}})\approx 0.0790.
\]

We can thus see, as discussed in class, BIC is more conservative since it selects the too-small AR(1) model fairly often, which shifts the post-selection distribution of \(\hat\rho_1\) to the right and makes it more dispersed. AIC selects the true AR(2) model more often, but it also over-selects AR(3) with nontrivial probability. In that sense, the Monte Carlo experiment shows very directly that once lag order is data-dependent, the distribution of the final estimator is a mixture over models rather than the distribution associated with a fixed AR(2) regression.

\section{Problem 3}

\paragraph{(i)}
Consider a univariate MA(1) process
\[
y_t = u_t + \theta u_{t-1},
\qquad
u_t \sim WN(0,\sigma_u^2).
\]
Its autocovariances are
\[
\Gamma_y(0)=\sigma_u^2(1+\theta^2),\qquad
\Gamma_y(1)=\sigma_u^2\theta,\qquad
\Gamma_y(\ell)=0 \ \text{for all } |\ell|>1.
\]
Thus, its spectral density is
\begin{align*}
s_y(\omega)
&= \frac{1}{2\pi}\sum_{\ell=-\infty}^{\infty}\Gamma_y(\ell)e^{-i\omega \ell} \\
&= \frac{1}{2\pi}\Bigl(\Gamma_y(0)+\Gamma_y(1)e^{-i\omega}+\Gamma_y(-1)e^{i\omega}\Bigr).
\end{align*}
Since the process is univariate and covariance stationary, we have
\[
\Gamma_y(-1)=\Gamma_y(1),
\]
so
\[
s_y(\omega)=\frac{1}{2\pi}\Bigl(\Gamma_y(0)+2\Gamma_y(1)\cos\omega\Bigr).
\]
Now a spectral density must satisfy
\[
s_y(\omega)\ge 0 \qquad \text{for all } \omega\in[-\pi,\pi].
\]
Therefore,
\[
\Gamma_y(0)+2\Gamma_y(1)\cos\omega \ge 0
\qquad \text{for all } \omega\in[-\pi,\pi].
\]
Since \(\cos\omega\in[-1,1]\), the smallest possible value of
\(\Gamma_y(0)+2\Gamma_y(1)\cos\omega\) is
\[
\Gamma_y(0)-2|\Gamma_y(1)|.
\]
Thus nonnegativity of the spectral density implies
\[
\Gamma_y(0)-2|\Gamma_y(1)|\ge 0,
\]
or equivalently,
\[
{\,|\Gamma_y(1)|\le \frac{\Gamma_y(0)}{2}\, .}
\]
Thus, there does not exist any univariate MA(1) process with
\[
|\Gamma_y(1)|>\frac{\Gamma_y(0)}{2}.
\]

As a check, using the MA(1) formulas above,
\[
\frac{|\Gamma_y(1)|}{\Gamma_y(0)}
=
\frac{|\theta|}{1+\theta^2}
\le \frac12,
\]
with equality if and only if \(|\theta|=1\), which is consistent with the spectral argument.

\paragraph{(ii)}
There is no contradiction with the Wold decomposition theorem.

The Wold theorem says that any covariance stationary process can be written as the sum of a deterministic component and a purely nondeterministic component, where the latter has a moving-average representation of possibly infinite order:
\[
y_t = d_t + \sum_{j=0}^{\infty}\psi_j \varepsilon_{t-j},
\qquad
\sum_{j=0}^{\infty}\psi_j^2<\infty,
\]
for a white noise process \(\{\varepsilon_t\}\).

Thus, Wold does \emph{not} say that every covariance stationary process is an MA(1). It only says that every covariance stationary process admits an MA(\(\infty\)) representation (up to a deterministic component). Therefore, part (i) only shows that the class of univariate MA(1) processes is too restrictive to generate autocovariances with
\[
|\Gamma_y(1)|>\frac{\Gamma_y(0)}{2}.
\]
A covariance stationary process with such autocovariances may still exist; it simply cannot be represented as an MA(1). In general it would require a longer, typically infinite-order, moving-average representation.

Therefore the ``paradox'' is resolved by noting that
\[
{\text{Wold } \neq \text{ MA(1)}.}
\]
Wold guarantees an MA(\(\infty\)) representation, not a finite MA(1) representation.


\section{Problem 4}

Let
\[
y_t^\dagger \equiv E^*(y_t \mid \{x_\tau\}_{-\infty<\tau<\infty})=\Theta(L)x_t,
\qquad
\Theta(z)\equiv \sum_{\ell=-\infty}^{\infty}\Theta_\ell z^\ell,
\]
where \(\Theta(z)\) is absolutely summable. Define the projection error
\[
e_t \equiv y_t-y_t^\dagger.
\]
Since \(y_t^\dagger\) is the best linear predictor of \(y_t\) from all leads and lags of \(\{x_t\}\), the projection theorem implies that \(e_t\) is orthogonal to the linear span of \(\{x_\tau\}_{-\infty<\tau<\infty}\). In particular,
\[
E(e_t x_{t-h}')=0 \qquad \text{for all } h\in\mathbb Z.
\]
Equivalently,
\[
\Gamma_{ex}(h)=0 \qquad \text{for all } h\in\mathbb Z.
\]

\paragraph{(i)}
Since
\[
y_t = y_t^\dagger + e_t,
\]
we have
\[
\Gamma_{yx}(h)=\Gamma_{y^\dagger x}(h)+\Gamma_{ex}(h)=\Gamma_{y^\dagger x}(h)
\qquad \text{for all } h\in\mathbb Z.
\]
Now
\[
y_t^\dagger=\sum_{\ell=-\infty}^{\infty}\Theta_\ell x_{t-\ell},
\]
so
\begin{align*}
\Gamma_{y^\dagger x}(h)
&= Cov\!\left(\sum_{\ell=-\infty}^{\infty}\Theta_\ell x_{t-\ell},\,x_{t-h}\right) \\
&= \sum_{\ell=-\infty}^{\infty}\Theta_\ell Cov(x_{t-\ell},x_{t-h}) \\
&= \sum_{\ell=-\infty}^{\infty}\Theta_\ell \Gamma_x(h-\ell).
\end{align*}

Taking Fourier transforms,
\begin{align*}
s_{yx}(\omega)
&= s_{y^\dagger x}(\omega) \\
&= \frac{1}{2\pi}\sum_{h=-\infty}^{\infty} e^{-i\omega h}\Gamma_{y^\dagger x}(h) \\
&= \frac{1}{2\pi}\sum_{h=-\infty}^{\infty} e^{-i\omega h}
\sum_{\ell=-\infty}^{\infty}\Theta_\ell \Gamma_x(h-\ell).
\end{align*}
Using absolute summability to interchange the sums, and setting \(m=h-\ell\), we get
\begin{align*}
s_{yx}(\omega)
&= \sum_{\ell=-\infty}^{\infty}\Theta_\ell e^{-i\omega \ell}
\left(\frac{1}{2\pi}\sum_{m=-\infty}^{\infty} e^{-i\omega m}\Gamma_x(m)\right) \\
&= \Theta(e^{-i\omega})\,s_x(\omega).
\end{align*}
Since \(s_x(\omega)\) is nonsingular for every \(\omega\in[-\pi,\pi]\), it follows that
\[
{\Theta(e^{-i\omega})=s_{yx}(\omega)s_x(\omega)^{-1}}
\qquad \text{for all } \omega\in[-\pi,\pi].
\]

\paragraph{(ii)}
Since \(y_t^\dagger=\Theta(L)x_t\), the spectrum of \(y_t^\dagger\) is
\[
s_{y^\dagger}(\omega)
=
\Theta(e^{-i\omega})\,s_x(\omega)\,\Theta(e^{-i\omega})^*.
\]
Using the result from part (i),
\[
\Theta(e^{-i\omega})=s_{yx}(\omega)s_x(\omega)^{-1}.
\]
Taking conjugate transposes and using that \(s_x(\omega)\) is Hermitian,
\[
\Theta(e^{-i\omega})^*
=\bigl(s_x(\omega)^{-1}\bigr)^* s_{yx}(\omega)^*
=s_x(\omega)^{-1}s_{xy}(\omega),
\]
because \(s_x(\omega)^*=s_x(\omega)\) implies \(\bigl(s_x(\omega)^{-1}\bigr)^*=s_x(\omega)^{-1}\), and
\[
s_{xy}(\omega)=s_{yx}(\omega)^*.
\]
Thus,
\begin{align*}
s_{y^\dagger}(\omega)
&= \Theta(e^{-i\omega})\,s_x(\omega)\,\Theta(e^{-i\omega})^* \\
&= s_{yx}(\omega)s_x(\omega)^{-1}\,s_x(\omega)\,
s_x(\omega)^{-1} s_{xy}(\omega) \\
&= s_{yx}(\omega)s_x(\omega)^{-1}s_{xy}(\omega).
\end{align*}
Therefore,
\[
{
s_{y^\dagger}(\omega)=s_{yx}(\omega)s_x(\omega)^{-1}s_{xy}(\omega)
}
\qquad \text{for all } \omega\in[-\pi,\pi].
\]

\paragraph{(iii)}
Now suppose \(n_y=n_x=1\), so all spectra are scalar-valued. The coherency is defined by
\[
K_{yx}(\omega)\equiv \frac{s_{yx}(\omega)}{\sqrt{s_y(\omega)s_x(\omega)}}.
\]
Since \(s_x(\omega)\) is nonsingular and scalar, we have \(s_x(\omega)>0\), and by assumption \(s_y(\omega)>0\). From part (ii),
\[
s_{y^\dagger}(\omega)=s_{yx}(\omega)s_x(\omega)^{-1}s_{xy}(\omega).
\]
In the scalar case,
\[
s_{xy}(\omega)=\overline{s_{yx}(\omega)},
\]
so
\[
s_{y^\dagger}(\omega)=\frac{|s_{yx}(\omega)|^2}{s_x(\omega)}.
\]
Dividing by \(s_y(\omega)\),
\[
\frac{s_{y^\dagger}(\omega)}{s_y(\omega)}
=
\frac{|s_{yx}(\omega)|^2}{s_y(\omega)s_x(\omega)}.
\]
But
\[
|K_{yx}(\omega)|^2
=
\left|\frac{s_{yx}(\omega)}{\sqrt{s_y(\omega)s_x(\omega)}}\right|^2
=
\frac{|s_{yx}(\omega)|^2}{s_y(\omega)s_x(\omega)}.
\]
Thus,
\[
{
|K_{yx}(\omega)|^2=\frac{s_{y^\dagger}(\omega)}{s_y(\omega)}
}
\qquad \text{for all } \omega\in[-\pi,\pi].
\]

In summary,
\[
{\Theta(e^{-i\omega})=s_{yx}(\omega)s_x(\omega)^{-1}},
\]
\[
{s_{y^\dagger}(\omega)=s_{yx}(\omega)s_x(\omega)^{-1}s_{xy}(\omega)},
\]
and, in the scalar case,
\[
{|K_{yx}(\omega)|^2=\frac{s_{y^\dagger}(\omega)}{s_y(\omega)}}.
\]
This shows that squared coherency is the frequency-by-frequency fraction of the spectrum of \(y_t\) explained by the best linear predictor based on all leads and lags of \(x_t\).

\section{Problem 5}

I used the monthly FRED unemployment-rate series \texttt{UNRATE} and \texttt{UNRATENSA}, restricted to 1948:1--2019:12. I compare a parametric AR spectral estimator with a smoothed-periodogram estimator to assess how much seasonal power remains after adjustment, especially near the annual frequency \(\omega=\pi/6\).

\paragraph{(i)}
For each series I estimate an AR($p$) with intercept by least squares and choose \(p\in\{1,\dots,50\}\) by BIC. The implied spectral estimator is
\[
\widehat{s}_y^{(p)}(\omega)
= \frac{\widehat{\sigma}^2}{2\pi \left|1-\sum_{\ell=1}^{p}\widehat{A}_{\ell}e^{-i\omega \ell}\right|^2},
\qquad \omega \in [0,\pi].
\]
I report the estimated log spectrum together with pointwise confidence bands based on delta-method standard errors, using the usual large-sample approximation that \(\widehat{\sigma}^2\) is asymptotically independent of the estimated AR coefficients. BIC selects \(p=\QFiveLagNSA\) for the unadjusted series and \(p=\QFiveLagSA\) for the adjusted series. The resulting AR log spectra are shown in Figure~3.

\begin{figure}[htbp]
\centering
\includegraphics[width=\textwidth]{\QFiveARFigurePath}
\caption{AR(BIC) log spectral density estimates with pointwise 95\% confidence bands.}
\end{figure}

\paragraph{(ii)}
I also estimate the spectrum nonparametrically by smoothing the periodogram with the Epanechnikov kernel
\[
w(x)=1-x^2,\qquad |x|\leq 1,
\]
using bandwidth \(B=\QFiveKernelBandwidth\) ordinates. This provides a less model-dependent check on the same spectral features. The estimator is
\[
\widehat{s}_y(\omega_j)
= \frac{1}{2\pi}\sum_{|k|\leq B} w_B(k)\,\widehat{I}(\omega_{j+k}),
\qquad
w_B(k)=\frac{w(k/B)}{\sum_{|m|\leq B} w(m/B)}.
\]
For the log spectrum I use the asymptotic normal approximation from the notes,
\[
\frac{\log \widehat{s}_y(\omega_j)-\log s_y(\omega_j)}
{\sqrt{\sum_{|k|\leq B} w_B(k)^2}}
\overset{a}{\sim} N(0,1),
\]
which yields a pointwise 95\% confidence band. The smoothed log spectra are shown in Figure~4.

\begin{figure}[htbp]
\centering
\includegraphics[width=\textwidth]{\QFiveKernelFigurePath}
\caption{Kernel log spectral density estimates with pointwise 95\% confidence bands.}
\end{figure}

\paragraph{(iii)}
Both estimators deliver the same basic conclusion: the unadjusted series has a pronounced concentration of power around the annual seasonal frequency \(\omega=\pi/6\), whereas the adjusted series is much flatter in that region. The comparison at \(\omega=\pi/6\) is summarized below:

\begin{center}
\begin{tabular}{lccc}
\toprule
Series & BIC lag & AR log spectrum at \(\pi/6\) & Kernel log spectrum at \(\pi/6\) \\
\midrule
Not seasonally adjusted & \QFiveLagNSA & \QFiveARAnnualNSA & \QFiveKernelAnnualNSA \\
Seasonally adjusted & \QFiveLagSA & \QFiveARAnnualSA & \QFiveKernelAnnualSA \\
\bottomrule
\end{tabular}
\end{center}

At \(\omega=\pi/6\), the unadjusted series has \(\QFiveARAnnualGap\) more log spectral power under the AR estimator and \(\QFiveKernelAnnualGap\) more under the kernel estimator. In both figures this appears as a clear hump at the 12-month cycle for \texttt{UNRATENSA} that is largely absent for \texttt{UNRATE}. On that basis, the seasonal adjustment appears to remove the annual component effectively while leaving the lower-frequency business-cycle variation intact.

\section{Problem 6}

For the NKPC estimation I use the seasonally adjusted FRED series \texttt{GDPDEF} and \texttt{UNRATE}. Inflation is measured by quarter-over-quarter log growth of the GDP deflator,
\[
\pi_t = \Delta \log(\text{GDPDEF}_t),
\]
and the slack variable \(x_t\) is the quarterly average unemployment rate. The largest overlapping quarterly sample before 2020 is \QSixRawStart--\QSixRawEnd\ (\QSixRawN\ quarters). After constructing the required leads and lags, the effective estimation sample becomes \QSixEstStart--\QSixEstEnd\ (\QSixEstN\ observations).

\paragraph{(i)}
Imposing rational expectations, I estimate the hybrid NKPC through the moment conditions
\[
E\!\left[\bigl(\pi_t-c-\gamma_f\pi_{t+1}-\gamma_b\pi_{t-1}-\lambda x_t\bigr)Z_{t-1}\right]=0,
\]
where
\[
Z_{t-1}=\bigl(1,\pi_{t-1},\pi_{t-2},\pi_{t-3},x_{t-1},x_{t-2},x_{t-3}\bigr)'.
\]
Writing
\[
u_t(\theta)=\pi_t-c-\gamma_f\pi_{t+1}-\gamma_b\pi_{t-1}-\lambda x_t,
\qquad
g_t(\theta)=Z_{t-1}u_t(\theta),
\]
the GMM estimator solves
\[
\widehat{\theta}(W)=\arg\min_{\theta}\,\bar g(\theta)'W\bar g(\theta),
\qquad
\bar g(\theta)=\frac{1}{T}\sum_{t=1}^{T} g_t(\theta).
\]
Using the weight matrix
\[
\widehat W_1=\left(\frac{Z'Z}{T}\right)^{-1}
\]
gives the 2SLS estimator requested in the question. Because the moments are linear in the parameters, the estimator has the closed form
\[
\widehat\theta(\widehat W)
=
\left(\frac{X'Z}{T}\widehat W\frac{Z'X}{T}\right)^{-1}
\left(\frac{X'Z}{T}\widehat W\frac{Z'y}{T}\right).
\]

\paragraph{(ii)}
I then compute the efficient two-step GMM estimator using
\[
\widehat W_2=\widehat\Omega(\widehat\theta_1)^{-1},
\]
where \(\widehat\Omega(\theta)\) is a Bartlett/Newey-West HAC estimator of the long-run variance of \(g_t(\theta)\):
\[
\widehat\Omega(\theta)
=
\widehat\Gamma_0
+
\sum_{\ell=1}^{L_T}
\left(1-\frac{\ell}{L_T+1}\right)\bigl(\widehat\Gamma_\ell+\widehat\Gamma_\ell'\bigr),
\]
\[
\widehat\Gamma_\ell
=
\frac{1}{T}\sum_{t=\ell+1}^{T}\tilde g_t(\theta)\tilde g_{t-\ell}(\theta)',
\qquad
\tilde g_t(\theta)=g_t(\theta)-\bar g(\theta).
\]
I use the automatic lag-truncation rule
\[
L_T=\left\lfloor 4(T/100)^{2/9}\right\rfloor=\QSixBandwidth,
\]
which is a standard weak-dependence consistent choice.

\paragraph{(iii)}
For both estimators I report HAC-robust standard errors from the usual GMM sandwich formula
\[
\widehat{\mathrm{Var}}(\widehat\theta)
=
\frac{1}{T}
(D'\widehat W D)^{-1}
\bigl(D'\widehat W\widehat\Omega \widehat W D\bigr)
(D'\widehat W D)^{-1},
\qquad
D=-\frac{Z'X}{T}.
\]
The resulting estimates are:

\begin{center}
\begin{tabular}{lcccc}
\toprule
Parameter & 2SLS / one-step GMM & s.e. & Efficient GMM & s.e. \\
\midrule
\(c\) & \QSixCOne & \QSixCOneSE & \QSixCTwo & \QSixCTwoSE \\
\(\gamma_f\) & \QSixGFOne & \QSixGFOneSE & \QSixGFTwo & \QSixGFTwoSE \\
\(\gamma_b\) & \QSixGBOne & \QSixGBOneSE & \QSixGBTwo & \QSixGBTwoSE \\
\(\lambda\) & \QSixLamOne & \QSixLamOneSE & \QSixLamTwo & \QSixLamTwoSE \\
\bottomrule
\end{tabular}
\end{center}

The estimates are stable across the two weighting matrices. In both specifications the forward-looking coefficient is large, around \(0.84\), the backward-looking coefficient is around \(0.19\) to \(0.21\), and the slope coefficient on unemployment is extremely small and statistically insignificant.

\paragraph{(iv)}
Using the efficient two-step estimator, Hansen's over-identification statistic is
\[
J
=
T\,\bar g(\widehat\theta_2)'\widehat\Omega(\widehat\theta_2)^{-1}\bar g(\widehat\theta_2).
\]
Since there are \(r=7\) instruments and \(k=4\) parameters, the null distribution is asymptotically \(\chi^2_{r-k}=\chi^2_3\). In this sample,
\[
J=\QSixJStat,
\qquad
p\text{-value}=\QSixJPValue.
\]
Thus the over-identifying restrictions are not rejected. Taken together with the parameter estimates, the evidence here points to a Phillips curve that is primarily forward-looking, with some backward-looking persistence, but with very little slope when slack is measured by the unemployment rate.

\end{document}
